\textbf{Enfin le cinquième et dernier ébranlement depuis 1950,}

Moins violente mais peut être tout aussi pernicieuse, en un certain sens, fut l’irruption de la
modernité : mécanisation de l’agriculture, arrivée des tracteurs, multiplication des moyens de
déplacement, ouverture sur le monde extérieur avec la radio puis la télévision vers 1960 ou
par les 130 harkis présents à Roquedols de 1962 à 1965 mais surtout l’essor du
tourisme « panacée perturbatrice » ?

On pourrait même parler de marée touristique surtout après la création du « Parc des
Cévennes » en 1970, malgré les nombreuses oppositions de ceux qui voyaient en lui la
création d’une « réserve cévenole » à l’image des « réserves indiennes » des Etats-Unis. Ne
disait-on pas que \emph{« les parcs naturels doivent être une contrepartie de la super civilisation technique de notre
ère : ils doivent jouer le rôle de sanatorium pour nos contemporains tourmentés dans les bureaux, les fabriques
et les usines... »}, écrivait Jean Fourastié.

Mais qu’on le redoute ou qu’on l’espère le Parc des Cévennes est l’avenir obligé du pays de
Meyrueis. \emph{« Parce que là seulement, avec les communautés de communes, peuvent se prendre des décisions
collectives, écrit Yves Grellier de Cabrillac en 1991. Les villages et les communes sont trop faibles pour
demeurer des instances délibératives, des lieux de décision. »}

Les conséquences sont multiples :

\begin{itemize}
    \item L’accélération de l’exode rural. Les petites exploitations disparaissent. Celles qui
    restent se mécanisent et s’agrandissent.

    \item Les bœufs font place aux tracteurs -- la dernière paire celle d’Emilien Martin de la
    Bragoussette en 1980- puis on voit apparaître les botteleuses pour le fourrage, les
    moissonneuses-batteuses... alors que disparaissent les troupeaux de vaches laitières,
    souvent remplacées par des chevaux de selle pour touristes voire pour émiris arabes.

    \item Les troupeaux de brebis se développent dans les vallées mais surtout sur les causses
    Méjean et Noir.

    \item Le Parc contribue à la réintroduction des vautours fauves, des cerfs, des biches et des
    chevreuils, des coqs de bruyère ou des castors. Mais l’emploi des pesticides et des
    engrais chimiques qui remplacent ou complètent le « migou » (fumier de brebis)
    ajouté à l’irrigation et aux adductions d’eau, a fait baisser le niveau des rivières et
    contribué à la disparition d’une bonne partie des truites.

    \item Les fermes sont modernisées et aménagées en gîtes ou appartements modernes pour
    recevoir des familles «de la ville».Une quarantaine à Meyrueis, dans les vallées ou sur
    le Causse, avec tout le confort, (de 200 à 400€ la semaine).

    \item Sept ou huit campings sont complets pendant les deux mois d’été, tout comme les six
    hôtels.

    \item Les produits des fermes ne sont plus transformées sur place : plus de lait, plus de
    tomes de chèvres, ou de plaquettes de beurre dans sa feuille d’oseille. mais des
    exploitations semi industrielles : le Fédou à Hylzas.

    \item On se recycle, en organisant des promenades à cheval, des randonnées en VTT, des
    écoles d’escalade dans les couronnes, on guide des groupes à la découverte des
    vautours, des grottes et avens, des cerfs à l’époque du brame, ou du massif de
    l’Aigoual...
\end{itemize}

Un monde s’est éteint. De nouveaux modes de vie sont apparus, le tourisme a remplacé les
modes de vie traditionnelle et Meyrueis avec sa maison de retraite et son foyer de vie pour
handicapés, sa fabrique de brosses pour « matériel de maintenance des surfaces », comme on
dit aujourd’hui, et son atelier de l’ONF fabriquant des mobiliers urbains ou routiers comme
des meubles de jardin, voit augmenter sa population.
